\documentclass[journal=jpcbfk, manuscript=article, layout=twocolumn, doi=true]{achemso}

\usepackage[utf8]{inputenc}
\usepackage{amsmath}
\usepackage{amssymb}
\usepackage{bm}
\usepackage{graphicx}
\usepackage{physics}
\usepackage{booktabs}
\usepackage{xcolor}
\usepackage{subcaption}

\makeatletter
\setlength{\@fptop}{0pt}
\setlength{\@dblfptop}{0pt}
\makeatother

\title{Site-Energy Detuning Explains the Apparent Excitonic Coupling of the
Venus\textsubscript{A206} Tandem Dimer}

\author{Robson Christie}
\affiliation{School of Mathematics and Physics, University of Portsmouth, Portsmouth PO1 3HF, United Kingdom}

\author{Cerys Murray}
\affiliation{School of Mathematics and Physics, University of Surrey, Guildford GU2 7XH, UK}

\author{Youngchan Kim}
\affiliation{Leverhulme Quantum Biology Doctoral Training Centre, School of Biosciences, Advanced Technology Institute, University of Surrey, Guildford GU2 7XH, UK}

\author{Jaewoo Joo}
\email{jaewoo.joo@port.ac.uk}
\affiliation{School of Mathematics and Physics, University of Portsmouth, Portsmouth PO1 3HF, United Kingdom}

\begin{document}

\begin{abstract}
Circular-dichroism (CD) analyses of the Venus\textsubscript{A206} tandem dimer
report excitonic couplings of $131$--$186~\mathrm{cm^{-1}}$, an order of
magnitude larger than the Coulomb coupling obtained from integrated transition
densities. We show that this discrepancy is not an error in the computed
coupling but a consequence of the degeneracy assumption embedded in the
experimental analysis. Domain-based local-pair-natural-orbital
similarity-transformed equation-of-motion coupled-cluster transition densities
integrated over $1000$ snapshots of the covalently tethered dimer give
$J = 32.82 \pm 1.55~\mathrm{cm^{-1}}$, validated against an independent
gas-phase control to $1\%$. Quantum-mechanics/molecular-mechanics calculations
on the same construct show that the two chromophores are \emph{not}
isoenergetic: the site-energy detuning is
$\langle|\Delta|\rangle = 603~\mathrm{cm^{-1}}$, an order of magnitude above
$2|J| = 65.6~\mathrm{cm^{-1}}$. Because the exciton-chirality rotational
strength of each band scales as $J/\Omega$ with
$\Omega = \sqrt{\Delta^{2}+4J^{2}}$, a CD measurement reports a
rotational-strength-weighted mean of $\Omega$ that is close to its harmonic
mean and roughly half its arithmetic mean. Interpreted as $2J$ under an
assumption of degeneracy, this yields an apparent coupling of
$105$--$159~\mathrm{cm^{-1}}$ depending on the treatment of environmental
screening, encompassing the constrained experimental value while the true
coupling remains $32.8~\mathrm{cm^{-1}}$. The ratio of the reported to the computed coupling,
$130.79/32.82 = 3.99$, is thus accounted for. The eigenstates are
correspondingly localised (mean mixing $0.18$), so the CD splitting of this
system should not be read as a Davydov splitting.
\end{abstract}

\section{Introduction}

Excitonic coupling in fluorescent-protein dimers has been inferred
spectroscopically from the Davydov structure of difference-CD
spectra.\cite{kim2019venusa206,nguyen2025anomalous} For the
Venus\textsubscript{A206} tandem dimer (dVenus-TD), a constrained
three-Lorentzian analysis places the excitonic pair
$261.58~\mathrm{cm^{-1}}$ apart, giving an apparent coupling
$U = 130.79~\mathrm{cm^{-1}}$; an unconstrained fit gives
$U = 185.4~\mathrm{cm^{-1}}$.\cite{nguyen2025anomalous} Couplings computed
from first principles by integrating monomer transition densities are
substantially smaller, and the resulting mismatch has been treated as a
deficiency of the electronic-structure model.

Both experimental analyses assume that the two chromophores are isoenergetic,
and in the constrained fit the assumption is not implicit but structural. That
model places the two couplet components at $G_{1} \mp \delta$, where $G_{1}$ is
held fixed at the first Gaussian of the corresponding one-chromophore
(dVenus-TDX) \emph{absorption} spectrum and $\delta$ is the only free splitting
parameter, so $\delta$ is the coupling by construction: a symmetric splitting
about the monomer energy is the degenerate-dimer prediction, and the detuning
$\Delta$ never enters as a parameter. The supplement gives the rationale as the
symmetry ``expected for Davydov splitting'' about that peak on an energy scale.
Two features of the same analysis make the assumption worth testing rather than
adopting. The centre is taken from the one-chromophore control, yet the
authors' own absorption decomposition reports that this component moves
$35~\mathrm{cm^{-1}}$ to the red --- $27\%$ of $\delta$ --- in the
two-chromophore construct. And the splitting recovered, $2\delta =
261.58~\mathrm{cm^{-1}}$, is smaller than the fitted half-width of either band
it separates ($293.5$ and $439.3~\mathrm{cm^{-1}}$), so the couplet is
unresolved and the fit is correspondingly ill-conditioned. Nguyen \emph{et al.}
are themselves cautious: they describe their couplings as ``apparent'', note
that the fitted components may result from ``the merger of multiple unresolved
split components'', and caution that a model with ``only two excited states
might be overly simplistic''.\cite{nguyen2025anomalous} Here we test the
degeneracy assumption directly and show that it is the origin of the
discrepancy.

\section{Methods}

\subsection{Construct}

The simulated construct is the covalently tethered tandem dimer of Kim
\emph{et al.}\cite{kim2019venusa206} A single $491$-residue polypeptide carries
two complete Venus $\beta$-barrels joined by a $33$-residue linker
(\texttt{SGLRSENLYFQGPREFCRYPAQWRPLESRPRTT}, containing the TEV site
\texttt{ENLYFQG}); barrel~1 spans residues $1$--$229$, the linker $230$--$262$,
and barrel~2 $263$--$491$. The chromophores are the anionic CR2 residues at
positions $66$ and $328$. The A206 dimer interface is docked throughout, with
$40$ residues in inter-barrel contact; mapping to GFP numbering (shifted by
$-1$ through the chromophore condensation), Ala206, Leu207, Leu221 and Phe223
all appear in the contact set. This is the construct whose sequence Nguyen
\emph{et al.} inherit by citation for dVenus-TD, with one stated difference:
for bacterial expression they mutated the single linker cysteine to glycine to
suppress spontaneous inter-construct disulfide
formation.\cite{nguyen2025anomalous} The residue is Cys246, in the flexible
tether and $23.8$ and $42.9~\text{\AA}$ from the two chromophore centroids in
the built structure; a trajectory rebuilt with the substitution gives
site-energy statistics indistinguishable from the ensemble used here
(Mann--Whitney $p = 0.45$).

\subsection{Coulomb coupling}

Inter-chromophore couplings were obtained by integrating the cap-masked
DLPNO-STEOM-CCSD transition density over both sites of each snapshot, as
described previously. To be explicit about the scheme: this is the
transition-density-cube approach of Krueger, Scholes and
Fleming,\cite{krueger1998} in which $q_i$ and $q_j$ in the Coulomb double sum
are the transition densities integrated over individual grid \emph{voxels} of
the cap-masked cube, not the atom-centred transition charges of a TrEsp or
Renger-type fit.\cite{madjet2006} No atomic-charge fitting is performed at any
stage. Distances in the Coulomb double sum are evaluated in
\AA{}, so the reciprocal distance is converted to atomic units with
$1~\text{\AA}^{-1} = 0.529177\,a_{0}^{-1}$. Using the coordinate conversion
$1.8897$ in its place inflates every coupling by
$1/0.529177^{2} = 3.5711$; this error is now guarded by a regression test that
checks a closed-form two-charge reference and scans the source for the
signature.

\subsection{Site energies}

Site energies were computed with time-dependent density-functional theory in
the Tamm--Dancoff approximation (CAM-B3LYP/6-311+G**, TeraChem 1.97B) on
$44$-atom QM regions embedded in the complete
$69{,}412$-point-charge field of the solvated system. The region is specified
exactly as follows, since the count is easy to reproduce incorrectly:

\begin{itemize}
\item the complete anionic CR2 chromophore residue, \textbf{29 atoms}
      (all heavy atoms and their hydrogens, the residue as parameterised in
      the AMBER topology);
\item the $\pi$-stacked Tyr203 \textbf{phenol group only}, \textbf{12 atoms}
      --- CG, CD1, CD2, CE1, CE2, CZ, OH and the five hydrogens HD1, HD2, HE1,
      HE2, HH. The tyrosine C$_\beta$ and backbone are \emph{not} included;
      the cut is made at the C$_\gamma$--C$_\beta$ bond;
\item \textbf{3 link hydrogens} capping the three bonds severed at the QM/MM
      boundary.
\end{itemize}

The stoichiometry is C$_{19}$H$_{18}$N$_3$O$_4$ with a net charge of $-1$, and
a structure diagram with every atom labelled is given as
Fig.~S\ref{fig:qmregion} in the Supporting Information. A count of $45$ arises
if the tyrosine C$_\beta$ is retained instead of being replaced by a link
hydrogen.

Two embedding choices proved essential. First, both sites must see the
identical whole system: a $12~\text{\AA}$ shell gives the two supposedly
equivalent site models MM fields whose net charges differ by exactly
$1.000000000000\,e$ at every frame, whereas full-system embedding makes them
identical to $\sim\!5\times10^{-15}\,e$ by construction. Second, only MM atoms
covalently bonded across the QM boundary may be excluded. A geometric
$1.8~\text{\AA}$ cutoff additionally deletes hydrogen-bond partners --- the
His HD1 proton, Glu OE2, and an Arg or Gln proton --- and membership of that
set changes with thermal hydrogen-bond length, injecting large
frame-dependent errors. Radius-truncated shells are unusable for ensembles:
shell net charge is frame-unstable ($\pm2$ to $\pm4\,e$), and even on a
charge-symmetric frame a $20~\text{\AA}$ shell errs by $+450.7~\mathrm{cm^{-1}}$
against the full-system reference.

To separate the intrinsic electronic structure from the electrostatic
environment, both methods were also run in the gas phase on the identical
frozen $44$-atom anion --- same charge, multiplicity, method, basis and
geometry, with every protein and solvent point charge removed.
Table~\ref{tab:gasphase} collects the pair.

\begin{table}[t]
\centering
\caption{Bright-state vertical excitation of the $44$-atom model with and
without the protein/solvent point-charge field. Geometry, charge,
multiplicity, method and basis are identical within each row; only the
embedding differs.}
\label{tab:gasphase}
\begin{tabular}{lccc}
\toprule
Method & Gas phase & Embedded & Environmental shift \\
       & (cm$^{-1}$ / nm) & (cm$^{-1}$ / nm) & (cm$^{-1}$) \\
\midrule
TDA-wB97X-D3/6-311G** & $27\,324.5$ / $366.0$ & $27\,710.2$ / $360.9$ & $+385.7$ \\
DLPNO-STEOM-CCSD/def2-SVPD & \emph{[running]} & $19\,088.2$ / $523.9$ & \emph{[running]} \\
\bottomrule
\end{tabular}
\end{table}

The bright state is the same state in both environments: the gas-phase and
embedded transition dipoles are collinear to $\cos\theta = 0.9992$. The
environmental shift is therefore a genuine electrostatic effect on a fixed
electronic state, not a change of state ordering. Note that the comparison
isolates electrostatics only: the geometry is held frozen, so it does not
include the geometric relaxation the protein also imposes.

The bright $\pi\pi^{*}$ state was assigned by maximum oscillator strength;
frames whose runner-up margin fell below $3\times$ were discarded rather than
averaged. Excited-state phase is arbitrary per calculation, so transition
dipoles were oriented along the chromophore long axis (imidazolinone ring
centroid to phenolate oxygen, $|\!\cos| \ge 0.994$) before forming the
chirality triple product; without this the product changes sign from frame to
frame and the ensemble CD couplet averages to zero.

\subsection{Exciton model}

For each snapshot the two-site Hamiltonian
$H = \left(\begin{smallmatrix} E_{A} & J \\ J & E_{B}\end{smallmatrix}\right)$ was
diagonalised, giving $\Omega = \sqrt{\Delta^{2} + 4J^{2}}$ with
$\Delta = E_{A}-E_{B}$. Exciton-chirality rotational strengths were taken as
$R_{\pm} = \mp\pi\bar\nu\,c_{A}c_{B}\,\mathbf{R}_{AB}\cdot(\bm\mu_{A}\times\bm\mu_{B})$
with $c_{A}c_{B} = \pm J/\Omega$.

\section{Results}

\subsection{The coupling is small and independently validated}

Across $1000$ snapshots the transition-density coupling is
$J = 32.82 \pm 1.55~\mathrm{cm^{-1}}$, only $1.19$ times the point-dipole
value. A method-matched gas-phase control gives distributed-density and
point-dipole magnitudes of $6.09$ and $5.36~\mathrm{meV}$, against
$6.02$ and $5.36~\mathrm{meV}$ from an independent Q-Chem calculation
performed by a referee: agreement to $1\%$ on the distributed density and to
three significant figures on the point-dipole term.

The excess of the distributed-density coupling over the point-dipole value is
accounted for term by term by a multipole decomposition of the same transition
density: dipole--dipole $24.81$, dipole--quadrupole $4.03$,
dipole--octupole $0.73$ and quadrupole--quadrupole $0.16~\mathrm{cm^{-1}}$,
summing to $29.74$ against the full $30.45~\mathrm{cm^{-1}}$. The near-field
correction at this separation is therefore a $19\%$ dipole--quadrupole effect,
not a factor of several.

An independent confirmation has since appeared. Cusick \emph{et
al.}\cite{cusick2026twophoton}, from the same laboratories that reported the
larger values, compute dipole--dipole couplings of $26.9$--$36.8~\mathrm{cm^{-1}}$
for the Venus dimer on the same 1MYW structure, adopting
$\mu = 7.9 \pm 0.5$~D and $R = 25.4 \pm 1.0$~\AA{} and no dielectric screening.
Repeating our own point-dipole evaluation on that structure under their
conventions gives $31.1~\mathrm{cm^{-1}}$ against their $32.1$, an agreement of
$3\%$ between independent implementations.

\subsection{Kim's Davydov splitting is $\Omega$, not $2|J|$}
\label{sec:omega}

The most direct test of the small coupling is the splitting Kim
\emph{et al.}\cite{kim2019venusa206} measured, $14.6 \pm 0.3$~nm, which has
been read throughout the literature as $2|J|$ and therefore as
$J = 274~\mathrm{cm^{-1}}$. That reading assumes degenerate sites. For a
detuned dimer the exciton gap is
\begin{equation}
\Omega = \sqrt{\Delta^{2} + 4J^{2}},
\label{eq:omega}
\end{equation}
which for the present system is dominated by the detuning rather than by the
coupling. With the computed $J = 32.82~\mathrm{cm^{-1}}$ and
$|\Delta| = 550~\mathrm{cm^{-1}}$, Eq.~\eqref{eq:omega} gives
$\Omega = 553.9~\mathrm{cm^{-1}}$, which at $516$~nm is
\textbf{$14.75$~nm} --- inside Kim's stated uncertainty, with no adjustable
parameter.

The agreement is robust to the two largest remaining systematics. Replacing
our transition dipole ($9.6$~D from the spectroscopically normalised STEOM
density) with the experimental $7.9$~D, and removing the optical screening
$\varepsilon_{\mathrm{opt}} = 1.77$, moves $J$ between $22.2$ and
$39.3~\mathrm{cm^{-1}}$ but moves the predicted splitting only from $14.69$ to
$14.79$~nm, because $\Delta \gg 2|J|$ makes $\Omega$ almost independent of $J$.

Reading $14.6$~nm as $2|J|$ instead requires $J = 274~\mathrm{cm^{-1}}$, which
at the crystallographic $25.4$~\AA{} separation demands a transition dipole of
$\approx 29$~D, or a separation of $\approx 10$~\AA{} --- the distance between
the two halves of a bacteriochlorophyll special pair. Neither is available to
two $\beta$-barrels of diameter $\geq 22$~\AA.

The corollary is that the couplet \emph{amplitude}, suppressed by
$2|J|/\Omega = 0.119$, is the observable that carries $J$; the couplet
\emph{width} is a lineshape observable that saturates at the bandwidth.

\subsection{The chromophores are strongly inequivalent}

\begin{table}[t]
\centering
\caption{Site-energy detuning of the tandem dimer from QM/MM
(CAM-B3LYP/6-311+G**, full-system embedding, link-only exclusion).}
\label{tab:detuning}
\begin{tabular}{lc}
\toprule
Quantity & Value (cm$^{-1}$) \\
\midrule
$\langle|\Delta|\rangle$              & $603.3$ \\
median $|\Delta|$                     & $429.2$ \\
sample standard deviation of $\Delta$ & $703.6$ \\
$\langle\Omega\rangle$ (arithmetic)   & $611.2$ \\
$\Omega$ (harmonic mean)              & $295.9$ \\
$2|J|$ for comparison                 & $65.63$ \\
mean mixing $2|J|/\Omega$             & $0.18$ \\
\bottomrule
\end{tabular}
\end{table}

Table~\ref{tab:detuning} collects the ensemble. The detuning exceeds $2|J|$ in
$32$ of $32$ frames at this level of theory and in $23$ of $25$ frames in a
cheaper basis; the mean detuning is $9.2$ times $2|J|$. The eigenstates are
therefore largely localised, with mean mixing $0.18$, corresponding to roughly
$98\%$ of the excitation residing on one chromophore at any instant.

The detuning decorrelates in $\le 5$~ps: nine frames spaced $5$~ps apart span
a standard deviation of $774.3~\mathrm{cm^{-1}}$, $1.10$ times that of the full
nanosecond ensemble, and the lag-one autocorrelation is indistinguishable from
zero. The $1$~ns trajectory therefore already contains $\sim\!200$ independent
samples, and no additional sampling is required.

\subsection{Why a CD analysis reports a larger coupling}

The rotational strength of snapshot $i$ is
$|R_i| \propto (J/\Omega_i)\,|T_i|$, where
$T_i = \mathbf{R}_{AB}\cdot(\bm\mu_{A}\times\bm\mu_{B})$ is the chirality
triple product. A CD measurement therefore reports
\begin{equation}
\langle\Omega\rangle_{\mathrm{CD}}
= \frac{\sum_i (J/\Omega_i)\,|T_i|\,\Omega_i}{\sum_i (J/\Omega_i)\,|T_i|} ,
\label{eq:weighted}
\end{equation}
which reduces to the harmonic mean $N/\sum_i \Omega_i^{-1}$ when $T_i$ is
constant. Because the $1/\Omega$ factor emphasises the smallest splittings,
rare near-degenerate snapshots carry disproportionate weight: three of the
present $32$ frames account for $31.4\%$ of the total CD weight and eight for
$60.1\%$.

Retaining the per-snapshot $|T_i|$ matters. The triple product varies by
$19.6\%$ across the ensemble and is positively correlated with the splitting
($r = +0.325$), partially offsetting the $1/\Omega$ weighting; the transition
dipole magnitudes themselves are nearly constant ($3.7\%$). Treating $T_i$ as
constant, i.e.\ using the plain harmonic mean, gives
$295.9~\mathrm{cm^{-1}}$ and understates the apparent coupling by $7.5\%$.
Equation~(\ref{eq:weighted}) gives $318.1~\mathrm{cm^{-1}}$, roughly half the
arithmetic mean of $611.2~\mathrm{cm^{-1}}$, corresponding to an apparent
coupling of $159.0~\mathrm{cm^{-1}}$.

\begin{table}[t]
\centering
\caption{Apparent coupling recovered by a degenerate-model analysis of the
detuned ensemble, as a function of the environmental screening applied to the
classical field, from Eq.~(\ref{eq:weighted}). Intervals are $16$--$84$
percentiles from $3000$ bootstrap resamples of the present $n=32$ ensemble.
The true coupling is $32.82~\mathrm{cm^{-1}}$ throughout.}
\label{tab:apparent}
\begin{tabular}{lccc}
\toprule
Screening & $\langle|\Delta|\rangle$ & $\langle\Omega\rangle_{\mathrm{CD}}$ & $J_{\mathrm{app}}$ \\
\midrule
none                           & $603.3$ & $318.1$ & $159.0^{+31}_{-23}$ \\
$\varepsilon_{\rm opt} = 1.77$ & $340.8$ & $210.3$ & $105.2^{+16}_{-13}$ \\
\midrule
\multicolumn{3}{l}{Exp.\ dVenus-TD, constrained} & $130.79$ \\
\multicolumn{3}{l}{Exp.\ dVenus-TD, unconstrained} & $185.4$ \\
\bottomrule
\end{tabular}
\end{table}

Table~\ref{tab:apparent} shows that the apparent coupling lies between $105$ and
$159~\mathrm{cm^{-1}}$ across the defensible range of environmental screening.
That interval contains the constrained experimental value of
$130.79~\mathrm{cm^{-1}}$ and falls $14\%$ short of the unconstrained
$185.4~\mathrm{cm^{-1}}$, which is the fit whose components are least
determined by the data (Sec.~\ref{sec:couplet}). In every case the apparent
coupling exceeds the true coupling by a factor of $3.2$--$4.8$. Taking the constrained
experimental value at face value,
\begin{equation}
\frac{U_{\mathrm{exp}}}{J_{\mathrm{calc}}}
= \frac{130.79}{32.82} = 3.99 ,
\end{equation}
so the long-standing ``factor of four'' is accounted for to within $1\%$ by the
degeneracy assumption alone.

\subsection{The couplet on an absolute scale}
\label{sec:absolute}

The interaction-induced couplet can be placed on an absolute scale. Nguyen
et al.\cite{nguyen2025anomalous} report
$\varepsilon_{516} = 184\,400~\mathrm{M^{-1}\,cm^{-1}}$ for dVenus-TD and
$92\,200$ for dVenus-TDX, all samples at $25~\mu\mathrm{M}$ in a $0.5$~cm
cuvette, with $[\theta] = (\mathrm{mdeg}\times M_{w})/(10\,L\,C)$. The exact
$2\!:\!1$ ratio of the two extinction coefficients fixes the normalisation as
per mole of construct, TD carrying two chromophores and TDX one, so the
Table~S3 Lorentzian amplitudes $A_{1} = +3.59$ and $A_{2} = -2.27$ (in
$10^{5}~\mathrm{deg\,cm^{2}\,dmol^{-1}}$, at $\bar\nu_{0}\mp130.79~
\mathrm{cm^{-1}}$) are absolute. Two conditions on that scale are checked by
the authors themselves. The chromophore count is verified independently of the
assumed $\varepsilon$: the TDX absorption peaks are $0.52$, $0.51$ and $0.54$
of their TD counterparts, bounding missing or immature chromophores in the
three TD samples at $3.4\%$, $2.7\%$ and $9.3\%$, so for dVenus-TD the
two-chromophore normalisation is good to $3\%$. Against that, all spectra were
recorded at $25~\mu\mathrm{M}$, where the mean intermolecular separation is
$\sim\!20$~nm, and the authors concede in their Discussion that even the TDX
negative control ``might have had some small degree of intermolecular excitonic
coupling'' at this concentration --- a contribution that would enter the
difference spectrum with the wrong sign.

No magnetic transition dipole is needed for the comparison, because the
coupled-oscillator couplet is generated by the two electric transition dipoles
and their displacement alone. The quantity to compare is the couplet first
moment,
\begin{equation}
M = \int \frac{\Delta\varepsilon(\bar\nu)}{\bar\nu}\,
        (\bar\nu-\bar\nu_{0})\,\mathrm{d}\bar\nu
  = \frac{R_{+}\Omega}{2.296\times10^{-39}}
  = -\,\frac{\pi\bar\nu_{0}\,J\,T}{2.296\times10^{-39}} ,
\label{eq:firstmomentcd}
\end{equation}
in which the $1/\Omega$ carried by $R_{\pm}$ cancels against the $\Omega$ lever
arm. $M$ therefore depends on the coupling and the chirality triple product
$T = \mathbf{R}_{AB}\cdot(\bm\mu_{A}\times\bm\mu_{B})$ alone, and is
independent of the detuning, of the homogeneous linewidth and of the
inhomogeneous spread---every quantity in this problem that is either uncertain
or force-field-derived.

Over the $95$ usable production frames $\langle T\rangle = -81.7 \pm
16.8~\mathrm{\AA}\,\mathrm{a.u.}^{2}$, and an independent $1000$-frame ensemble
built on transition-density rather than heavy-atom centroids gives $-69.6 \pm
11.7$, so the geometric systematic is $15\%$ and the sign is unanimous in both.
The experimental couplet must be extracted with care. Table~S3 fits the
\emph{difference} spectrum, which Nguyen \emph{et al.} define by ``subtracting
the TD from TDX'', i.e.\ $\Delta[\theta] = [\theta]_{\mathrm{TDX}} -
[\theta]_{\mathrm{TD}}$. The couplet exists only in TD, so the couplet carried
by $\Delta[\theta]$ is \emph{minus} that of the dimer itself: their tabulated
$A_{1} = +3.59$ sits on the low-wavenumber component, and inverting it gives a
dimer couplet that is positive to the blue and negative to the red, a negative
exciton chirality. Equation~(\ref{eq:firstmomentcd}) gives
$M_{\mathrm{calc}} = +1.5\times10^{7}$ against
$M_{\mathrm{exp}} = +3.7\times10^{6}~\mathrm{deg\,cm^{2}\,dmol^{-1}\,cm^{-1}}$,
so the handedness agrees and the magnitudes differ by a factor of $4.1$ for a
quantity with no adjustable scale.

Peak-to-peak intensities agree far better---$2.9\times10^{5}$ predicted against
$2.8\times10^{5}$ observed, dissymmetry factors $4.8\times10^{-4}$ and
$4.7\times10^{-4}$---but that agreement is partly fortuitous. Our couplet is
spread over a mean $\Omega$ of $584~\mathrm{cm^{-1}}$ against the observed
$262~\mathrm{cm^{-1}}$, and the peak amplitude lost to that extra width offsets
the fourfold excess in $M$. The first moment is the quantity to trust.

That fourfold excess is most plausibly the transition-dipole normalisation
rather than the geometry, because $M \propto J\,T$ with both $J$ and $T$
quadratic in $|\bm\mu|$, so $M \propto |\bm\mu|^{4}$ and the excess is only a
$40\%$ error in the dipole. The published extinction coefficient now bounds
that independently. Scaling the Table~S2 four-Gaussian decomposition of the
dVenus-TDX band to $\varepsilon_{516} = 92{,}200~\mathrm{M^{-1}cm^{-1}}$ and
integrating $\int(\varepsilon/\bar\nu)\,\mathrm{d}\bar\nu$ gives
$|\bm\mu| = 7.5$~D reading the tabulated widths in the convention of the
fitting package the authors name, and $6.2$--$8.6$~D across the plausible width
conventions, against the $9.8$~D carried by our spectroscopically normalised
STEOM density. A dipole at the top of that range removes most of the factor of
four; one at $7.3$~D removes all of it. We do not rescale, because an
extinction-derived dipole carries the medium local-field factor and would
double-count the explicit $\varepsilon_{\mathrm{opt}} = 1.77$ screening applied
to $J$, but the comparison shows that the absolute CD intensity does not
constrain the geometry independently of that normalisation.

\subsection{Two observables that constrain the coupling without assuming
degeneracy}

The CD splitting is not the only coupling-sensitive quantity Nguyen
\emph{et al.}\cite{nguyen2025anomalous} measured. Two others were reported but
not analysed as coupling measurements, and both are free of the degeneracy
assumption.

\emph{Radiative rate.} For the two-site Hamiltonian of the Methods the
eigenstate dipole strengths are
$|\bm{\mu}_\pm|^2 = \mu^2\left(1 \pm \sin 2\theta\,\cos\alpha\right)$ with
$\sin 2\theta = 2J/\Omega$ and $\alpha$ the angle between the monomer
transition dipoles, so the radiative rate of the emitting state is shifted from
the monomer value by exactly $(2J/\Omega)\cos\alpha$. Nguyen \emph{et al.}\
report that dVenus-TD decays $57 \pm 4$~ps faster than the single-chromophore
dVenus-TDX control on a $3.026$~ns lifetime, and attribute this to
superradiance. With the Venus fluorescence quantum yield $\Phi = 0.57$ this
fixes
\begin{equation}
\left| \frac{2J}{\Omega}\cos\alpha \right|
= \frac{1}{\Phi}\,\frac{\Delta\tau}{\tau} = 0.033 .
\label{eq:superradiance}
\end{equation}
Their time-resolved anisotropy independently fixes the geometric factor: the
limiting anisotropy falls from $0.52$ in TDX to $0.30$ in TD, which under their
own complete-transfer analysis gives $|\cos\alpha| = 0.660$. Together these
require a mixing of $2J/\Omega = 0.050$. At $J = 32.82~\mathrm{cm^{-1}}$ that
implies $|\Delta| = 1310~\mathrm{cm^{-1}}$, the same order as the
$603~\mathrm{cm^{-1}}$ computed here. At $J = 130.79~\mathrm{cm^{-1}}$ it
implies $|\Delta| = 5219~\mathrm{cm^{-1}}$, or $0.65$~eV, which is not a
physically available detuning between two copies of the same chromophore. A
degenerate pair at $J = 130.79~\mathrm{cm^{-1}}$ would give
$2J/\Omega = 1$ and a lifetime shortening of $\sim\!1.1$~ns, twenty times what
was measured. The measured superradiance therefore excludes the degenerate
strong-coupling interpretation on its own.

\emph{Absorption shift.} The dipole-strength-weighted first moment of the dimer
absorption is
\begin{equation}
\frac{|\bm{\mu}_+|^2 E_+ + |\bm{\mu}_-|^2 E_-}{|\bm{\mu}_+|^2 + |\bm{\mu}_-|^2}
= \bar{E} + \frac{\Omega}{2}\,\frac{2J}{\Omega}\cos\alpha
= \bar{E} + J\cos\alpha ,
\label{eq:firstmoment}
\end{equation}
because the $\Omega$ carried by the intensity asymmetry cancels the $\Omega$
carried by the splitting. The shift is therefore independent of the detuning
for every individual dimer, and survives ensemble averaging over any
distribution of $\Delta$. It is the one coupling observable that a detuned
dimer and a degenerate dimer report identically. Nguyen \emph{et al.}\ report
it directly: the narrowest Gaussian of the absorption decomposition, $G_{1}$,
moves $35~\mathrm{cm^{-1}}$ to the red between dVenus-TDX and dVenus-TD
($517.5 \rightarrow 518.5$~nm, their Table~S2), and Kim \emph{et al.}\ obtain
the same value on a different sample, $33.8~\mathrm{cm^{-1}}$, from the
monomer-to-dimer shift of free Venus\textsubscript{A206}
($515.9 \rightarrow 516.8$~nm).\cite{kim2019venusa206} Two independent
measurements on different species agree to $4\%$. Combined with
$|\cos\alpha| = 0.660$ this bounds $|J| \le 53~\mathrm{cm^{-1}}$, the
inequality holding because the environmental change on dimerisation contributes
to the shift as well; reconciling $J = 130.79~\mathrm{cm^{-1}}$ with a
$35~\mathrm{cm^{-1}}$ shift instead requires $|\cos\alpha| \le 0.27$, dipoles
within $16^{\circ}$ of perpendicular, which is the opposite of the
``predominant head-to-tail'' orientation the same red shift is cited as
evidence for. The same $G_{1}$ shifts for their other two constructs are
$9~\mathrm{cm^{-1}}$ (dLanYFP-TD) and $29~\mathrm{cm^{-1}}$ (dTomato-TD),
against CD-derived couplings of $101.74$ and $134.39~\mathrm{cm^{-1}}$: in all
three constructs the absorption shift is four to eleven times smaller than the
coupling inferred from the CD couplet, and does not track it.

\subsection{Vibronic structure}

The third band in the experimental decomposition, at $481.37$~nm
($20\,774~\mathrm{cm^{-1}}$), lies $1452~\mathrm{cm^{-1}}$ above the couplet
centre $G_{1}$ --- the C=C/C=N stretch. Nguyen \emph{et al.} fit it in both
their models but do not assign it, and it is not negligible: its integrated
area is $18\%$ of the positive couplet lobe, and it carries the same sign as
the negative lobe, so it overlaps and biases the couplet it flanks. A
one-particle Holstein model with $\omega = 1450~\mathrm{cm^{-1}}$ (fixed by
that spacing, not fitted) places $21$--$24\%$ of the intensity in the
$470$--$495$~nm window for Huang--Rhys factors $S = 0.35$--$0.50$, against
$9.6\%$ for a purely electronic model, and at $S = 0.5$ resolves it as a
distinct band. The third component is therefore a vibronic sideband, not a
third electronic state, consistent with the caution of Nguyen
\emph{et al.}\cite{nguyen2025anomalous} that their components may merge
unresolved structure. The assignment also fixes the integration window that a
first-moment reduction of the couplet (Eq.~(\ref{eq:firstmomentcd})) requires:
without it the measured band does not integrate to zero and its first moment is
origin-dependent, which is why that reduction must be applied to the raw
difference spectrum rather than to the published decomposition.

\section{Discussion}

\subsection{What the CD couplet can and cannot determine}
\label{sec:couplet}

The width of the CD couplet does \emph{not} discriminate between a degenerate
and a detuned dimer. At a common linewidth the degenerate model gives couplet
separations of $9.2$--$13.7$~nm and the detuned model $12.8$--$17.4$~nm across
the range of component widths reported experimentally
($\mathrm{HWHM} = 294$--$439~\mathrm{cm^{-1}}$), against a measured
$14.6\pm0.3$~nm.\cite{kim2019venusa206} That measurement is a peak-to-peak
separation and so is the model quantity, but it is made on free, concentration-
driven Venus\textsubscript{A206} dimers rather than on the tethered tandem; the
tandem couplet is reported instead as fitted component centres, $7.0$~nm
(constrained) and $9.9$~nm (unconstrained).\cite{nguyen2025anomalous} Varying
the linewidth alone changes the modelled couplet by a factor of $4.4$, whereas
switching detuning on changes it by $\sim\!40\%$. The couplet is thus
consistent with our model but is not evidence for it; the evidence for detuning
is the electronic-structure calculation, not the spectrum.

The reason the couplet is so uninformative is that it is unresolved, and the
published fits establish this from their own parameters. A bisignate pair
separated by $\Omega$ and broadened by a lineshape of width $w \gg \Omega$
collapses to the derivative of that lineshape, whose extrema sit at
$\pm\sigma$ for a Gaussian and at $\pm\gamma/\sqrt{3}$ for a Lorentzian,
independently of $\Omega$. In that limit the peak positions carry the linewidth
and only the couplet amplitude and its signed first moment
(Eq.~(\ref{eq:firstmomentcd})) carry the coupling. The dVenus-TD fit is in this
regime by its own numbers: $2\delta = 261.58~\mathrm{cm^{-1}}$ against
Lorentzian half-widths of $293.5$ and $439.3~\mathrm{cm^{-1}}$. The
corresponding floor on any extremum-based reading is
$2\gamma/\sqrt{3} = 339$--$507~\mathrm{cm^{-1}}$, i.e.\ $9.0$--$13.5$~nm at
$516$~nm --- so a couplet of \emph{zero} splitting with these widths would
already reproduce most of the observed separation. Fitting component centres,
as both published analyses do, escapes that floor but not the underlying
ill-conditioning: as $\delta/\gamma \to 0$ the difference spectrum depends on
the amplitude and $\delta$ only through their product, so $\delta$ is fixed
solely by the departure from the derivative limit. At
$2\delta/\gamma = 0.6$--$0.9$ that departure is present but small, which is
consistent with Table~S3 quoting $\delta$ to five significant figures without
an interval despite the $5000$-sample resampling procedure described alongside
it. The same analysis also gives the two couplet components independent widths
differing by a factor of $1.50$, so the common-lineshape condition under which
a component separation maps cleanly onto a splitting does not hold within the
fit either. This joint indeterminacy of $J$, $\Delta$ and linewidth is why a
coupling inferred from a normalised CD profile is not uniquely determined.

\subsection{Consequences for the interpretation of coherence}

With mean mixing $0.18$ the exciton eigenstates are close to localised, so the
description of the absorbing state as delocalised over both chromophores does
not survive at the computed detunings. The photon-antibunching
observation\cite{kim2019venusa206} is unaffected, since it requires a single
emitting entity rather than a delocalised one.

\subsection{Limitations}
\label{sec:limitations}

Three limitations bound the quantitative claims. First, the classical
non-polarisable force field overestimates electrostatic gap fluctuations: a
second-order cumulant built from the measured gap correlation function gives a
dephasing time of $24$~fs, against $128.0\pm8.1$~fs measured for
dVenus-TD,\cite{nguyen2025anomalous} implying that the fast component of the
fluctuation is overestimated by $4.34\times$. The comparison is indicative
rather than exact: the measured value comes from two-photon near-field
scanning optical microscopy on a vacuum-dried thin film (their Table~S4), not
from the buffered solution in which the CD and absorption were recorded. Applying that factor to the total
detuning would reduce the apparent coupling to $60~\mathrm{cm^{-1}}$; we do not
do so, because the calibration is derived from fast fluctuations while the
detuning also contains slow structural contributions that are not a dielectric
screening effect. Second, the site-energy difference is sensitive to basis set
at the $27\%$ level.

Third, the time dependence of the screening within the dynamics window is
immaterial to every spectroscopic conclusion. Over the sub-picosecond window of
the electronic dynamics the inverse dielectric function relaxes only from
$1/\varepsilon = 0.56$ to $0.50$, which moves $J$ from $32.5$ to
$29.0~\mathrm{cm^{-1}}$, a $10.7\%$ change, and lengthens the nominal Rabi
period $h/2|J|$ from $0.51$ to $0.57$~ps. It moves the exciton gap $\Omega$ by
only $0.14\%$, and the minor-site population weight from $0.69\%$ to $0.55\%$.
The reason is Eq.~\eqref{eq:omega}: with $|\Delta| \approx 550~\mathrm{cm^{-1}}$
against $2|J| \approx 65~\mathrm{cm^{-1}}$, the gap and therefore the splitting,
the mixing and the exciton composition are set by the detuning, and are almost
independent of the coupling. A $10\%$ drift in $J$ during the first picosecond
cannot change a system that is $99\%$ localised into one that is not.

Fourth, the absolute CD is limited, but not in the way we previously stated.
TeraChem does print magnetic transition dipoles and both length- and
velocity-gauge rotational strengths for every root; they are nonetheless
unusable here. Across the $220$ completed site calculations the bright-state
velocity transition dipole carries only $0.233$ times the magnitude required by
the off-diagonal hypervirial relation
$\bm\mu_{\mathrm{vel}} = -\omega\,\bm\mu_{\mathrm{len}}$, and lies
$129.4\pm6.8^{\circ}$ away from it rather than the $180^{\circ}$ that relation
demands. The two gauges consequently return oscillator strengths differing by a
factor of $19$ ($f_{\mathrm{len}} = 1.13$, $f_{\mathrm{vel}} = 0.062$), and
rotational strengths that disagree in sign in $208$ of the $220$ runs. The
bright state carries essentially all of the chromophore's absorption intensity,
so $f_{\mathrm{vel}} = 0.06$ is untenable and it is the velocity-gauge quantity
that is at fault.

This is a property of the method rather than of the code. An independent PySCF
calculation on the same geometry and functional, also in the Tamm--Dancoff
approximation, returns $f_{\mathrm{len}}/f_{\mathrm{vel}} = 19.1$ in 6-31G* and
$17.0$ in 6-311+G**, with the velocity dipole at $0.229$ and $0.243$ of its
exact value---reproducing the TeraChem ensemble to within its scatter, and not
improving when diffuse functions are added. Two approximations are implicated:
the Tamm--Dancoff truncation discards the de-excitation block that the
off-diagonal hypervirial relation requires, and the range-separated exchange
operator does not commute with $\mathbf{r}$. A full linear-response treatment
would be needed before either gauge could be trusted.

This is not a cosmetic problem for the length-gauge value
either: a shift of the gauge origin by $\mathbf{a}$ moves it by
$\mathbf{a}\cdot(\bm\mu_{\mathrm{vel}}\times\bm\mu_{\mathrm{len}})$, which
vanishes only when the two dipoles are collinear, so here a displacement of
$0.79$~\AA{} changes $R(\mathrm{len})$ by $100\%$. Gauge disagreement and
origin dependence are the same defect measured two ways, and either one
disqualifies these numbers.

The absolute comparison of Sec.~\ref{sec:absolute} avoids them by construction,
needing only electric transition dipoles and geometry. What remains genuinely
unavailable is the \emph{intrinsic} monomer rotational strength, so we compare
the interaction-induced couplet rather than the full CD spectrum. The
experimental fit bounds what that omits: the two couplet bands of Table~S3 have
areas that cancel to $5.4\%$, so any non-conservative monomer contribution
beneath the band is small. What limits the absolute intensity is instead the
transition-dipole normalisation, through the fourth-power dependence set out in
Sec.~\ref{sec:absolute}; the handedness agrees once Nguyen's difference
spectrum is inverted to recover the dimer's own couplet.

A fourth limitation concerns sampling. The detuning itself is well determined:
it decorrelates in $\le 5$~ps, so the trajectory supplies $\sim\!200$
independent samples and $\langle|\Delta|\rangle$ is stable. The
rotational-strength-weighted mean of Eq.~(\ref{eq:weighted}) converges more
slowly, however, because the $1/\Omega$ factor concentrates weight on rare
near-degenerate snapshots. Bootstrap resampling of the present $n=32$ ensemble
gives a $16$--$84$ percentile interval of $\pm17\%$ on $J_{\mathrm{app}}$,
narrowing to $\pm12\%$ at $n=64$ and $\pm6\%$ at $n=256$, with a median stable
at $159$--$163~\mathrm{cm^{-1}}$ throughout. The numbers reported here should
therefore be read with that interval attached; a $200$-frame ensemble on the
$5$~ps grid is in progress and will reduce it to $\sim\!7\%$.

A fifth concerns the dimer geometry, and it is the one point at which our
structure and the experiment disagree. Our ensemble places the two transition
dipoles at $100.8 \pm 2.7^\circ$ ($|\cos\alpha| = 0.187$); the crystal
A206 dimer from which the construct was built gives $110.4^\circ$, and a
trajectory rebuilt with ff19SB/OPC gives $95.8 \pm 3.3^\circ$, so the value is
robust to force field and is set by the A206 interface rather than by the
sampling. The limiting-anisotropy drop requires $|\cos\alpha| = 0.660$ under
the assumption of complete excitation transfer, a factor of $3.5$ larger. Kim
\emph{et al.} measured the same quantity on the same construct and obtained
$r_{0} = 0.20$ against a monomer reference of $0.42$, giving
$|\cos\alpha| = 0.549$; the two reports differ by $20\%$ on the tandem and by
a comparable amount on the single-chromophore reference, from which the
intrinsic absorption--emission angle of one chromophore comes out as
$25^{\circ}$ and $14^{\circ}$ respectively. All four cross-combinations of the
two numerators and two references lie between $0.55$ and $0.79$, so the
qualitative disagreement with our structure is robust even though the value is
not well determined.

We can exclude several origins for it. The angle is not an artefact of placing
a rigid density: the rigid STEOM placement ($100.8^{\circ}$), the per-frame
QM/MM TDDFT dipoles computed independently at each geometry
($96.8 \pm 3.2^{\circ}$) and the bare chromophore long axes with no dipoles at
all ($98.6 \pm 3.0^{\circ}$) agree. It is not the transition-dipole direction
within the chromophore: STEOM-CCSD and CAM-B3LYP agree to $3.4^{\circ}$ and
both lie within $3.5^{\circ}$ of the long axis, so the $\sim\!12^{\circ}$
intramolecular tilt that would reconcile them is absent at two levels of
theory. And it is not the choice of crystallographic dimer. The deposited
1MYW biological unit is an exact two-fold, for which
$\cos\alpha = 2\cos^{2}\theta - 1$ with $\theta$ the angle between a
chromophore dipole and the C2 axis; the observed $\theta = 53.3^{\circ}$
reproduces $\cos\alpha = -0.285$ to machine precision, and the anisotropy would
require $\theta = 65.6^{\circ}$. Applying all six space-group operators over
the surrounding lattice returns exactly one further contact set, at
chromophore separations of $36$ and $47~\text{\AA}$, so no alternative packing
register combines the $\sim\!25~\text{\AA}$ separation both experimental
analyses assume with a larger $|\cos\alpha|$.

What we cannot exclude is that the tethered tandem does not adopt the free
dimer's register. The crystal register requires the $33$-residue linker to
span $54.2~\text{\AA}$ from the C-terminus of barrel~1 to the N-terminus of
barrel~2, against $\sim\!38~\text{\AA}$ for an unperturbed disordered chain of
that length, and across all $1000$ production frames the linker sits at
$52.3 \pm 0.7~\text{\AA}$ --- taut, with no slack, for a nominally flexible
tether. The construct was built in that register and a nanosecond cannot
re-dock two $\beta$-barrels, so no alternative has been sampled. This is the
one remaining candidate and it requires enhanced sampling to test.

Through Eq.~(\ref{eq:firstmoment}) our geometry and coupling predict an
absorption shift of $6.1~\mathrm{cm^{-1}}$ against the $35.3~\mathrm{cm^{-1}}$
measured. That comparison is not, however, the test it appears to be, because
the measured shift is not purely excitonic. The TDX control deletes Tyr67, so
its partner barrel is still folded and docked but its chromophore is absent ---
and with it the $-1~e$ that the anionic chromophore carries $24.7$~\AA{} away.
We isolated that term by repeating the site calculation on $16$ frames with
everything held fixed except the $29$ partner-chromophore point charges, which
were made neutral: same QM geometry, same TeraChem input, same
$69{,}374$-charge field at the same coordinates. Over the $15$ frames passing
the bright-state margin criterion the partner charge red-shifts both sites
equally, by $15.8 \pm 7.3~\mathrm{cm^{-1}}$ at site~A and
$15.5 \pm 4.5~\mathrm{cm^{-1}}$ at site~B, giving a band shift of
$15.6 \pm 4.0~\mathrm{cm^{-1}}$. A classical estimate over all $100$ available
frames, using the computed difference dipoles
($|\Delta\bm\mu| = 3.7$ and $3.0$~D) and the unscreened monopole field, gives
$97 \pm 4~\mathrm{cm^{-1}}$, so the self-consistent embedding screens the term
by a factor of $6$. Electrostatics and coupling together predict
$21.7 \pm 4.0~\mathrm{cm^{-1}}$, leaving $13.6~\mathrm{cm^{-1}}$ unassigned.
The conclusion we draw is negative and deliberately so: a demonstrably
non-excitonic contribution accounts for $44\%$ of the measured shift, so the
absorption shift cannot be used to infer a coupling in either direction, and it
is not evidence for a coupling larger than $32.82~\mathrm{cm^{-1}}$. We note that the
same tension appears within the experimental data alone: for all three
constructs the anisotropy and the absorption shift require values of
$|\cos\alpha|$ that differ by factors of $3.6$, $3.9$ and $7.1$. The
superradiance constraint of Eq.~(\ref{eq:superradiance}) is unaffected, since
it fixes the product $(2J/\Omega)\cos\alpha$ directly.

None of these affects the central conclusion, which depends only on the
inequality $|\Delta| \gg 2|J|$; that inequality holds under every embedding,
basis, functional and screening treatment examined here, and in every
individual snapshot.

\section{Conclusions}

The excitonic coupling of the Venus\textsubscript{A206} tandem dimer is
$32.82 \pm 1.55~\mathrm{cm^{-1}}$, not the $131$--$186~\mathrm{cm^{-1}}$
inferred from circular dichroism. The two values are reconciled by the
site-energy detuning, $\langle|\Delta|\rangle = 603~\mathrm{cm^{-1}}$, which
the experimental analysis does not include: because circular dichroism weights
each configuration by $J/\Omega$, it reports the harmonic mean of the exciton
splitting, and reading that as $2J$ inflates the inferred coupling roughly
fourfold. The dimer is therefore an inhomogeneously detuned pair of largely
localised chromophores rather than a degenerate excitonic dimer, and its CD
splitting should not be interpreted as a Davydov splitting.

\begin{acknowledgement}
We thank the anonymous referee whose independent Q-Chem calculation
corroborated the corrected transition-density coupling.
\end{acknowledgement}

\bibliography{Bibliography}

\end{document}
